\chapter{Anomaly scoring}
\label{sec:scoring}

\section{Composición}

GraphHunter asigna a cada path resultado de un hunt un \emph{score}
compuesto en $[0, 1]$, calculado como combinación ponderada de cinco
componentes independientes.

\begin{equation*}
S(p) = \sum_{i=1}^{5} w_i \cdot C_i(p) \qquad \sum_i w_i = 1
\end{equation*}

Los pesos $w_i$ están en \code{ScoringWeights} (ver
\filepath{core/graph-engine/src/anomaly.rs}). Valores por defecto:

\begin{table}[H]
\centering
\small
\begin{tabular}{@{}lll@{}}
\toprule
Campo & Peso default & Componente \\
\midrule
\code{w1_entity_rarity}      & 0.25 & Entity Rarity \\
\code{w2_edge_rarity}        & 0.30 & Edge Rarity \\
\code{w3_neighborhood_conc}  & 0.25 & Neighborhood Concentration \\
\code{w4_temporal_novelty}   & 0.20 & Temporal Novelty \\
\code{w5_gnn_threat}         & 0.00 & GNN Threat \\
\bottomrule
\end{tabular}
\caption{Pesos por defecto. El GNN está desactivado salvo que el
usuario provea modelo y opte-in.}
\end{table}

Los pesos se renormalizan automáticamente si no suman 1. El
usuario puede sobreescribirlos por sesión vía
\code{POST /enable_anomaly_scoring} o \code{POST /compute_scores}.

\section{Componente 1: Entity Rarity}
\label{sec:scoring-er}

Mide qué tan infrecuentes son las entidades del path individualmente.

\begin{equation*}
ER(v) = 1 - \frac{\ln(\text{freq}(v))}{\ln(\text{max\_freq})}
\qquad C_1(p) = \frac{1}{|p|} \sum_{v \in p} ER(v)
\end{equation*}

Donde $\text{freq}(v)$ es el conteo de apariciones del nodo $v$ en
eventos ingresados, y $\text{max\_freq}$ el máximo global. Un nodo
que aparece una vez tiene $ER \approx 1$; el nodo más visto tiene
$ER = 0$.

\textbf{Uso típico}: captura \emph{qué tan raro es ver a estos
actores del todo}. Sirve para priorizar paths donde participan
cuentas o hosts rara vez observados.

\section{Componente 2: Edge Rarity}
\label{sec:scoring-edger}

Mide qué tan infrecuentes son las transiciones observadas.

\begin{equation*}
EdgeR(s, d) = 1 - \frac{\ln(\text{freq}(s, d))}{\ln(\text{max\_pair\_freq})}
\end{equation*}

$C_2(p)$ se calcula como promedio sobre todas las aristas del path.

\textbf{Uso típico}: captura \emph{qué tan inusual es esta secuencia
concreta}. Un usuario que rutinariamente se autentica en un host
tiene $ER$ bajo para ambos nodos por separado, pero si un día ese
usuario escribe en un directorio sensible al que nunca había ido, el
edge rarity se dispara.

\section{Componente 3: Neighborhood Concentration}
\label{sec:scoring-nc}

Promedio sobre los nodos del path de su concentración de vecindario
(precomputada en \code{gm_graph_finalize}).

\begin{equation*}
NC(v) = \frac{\max_t |\{u : (v, u) \in E, \tau(u) = t\}|}{|\text{neigh}(v)|}
\end{equation*}

Donde $\tau(u)$ es el tipo de la entidad $u$. Un nodo cuyos vecinos
son de tipos muy diversos tiene $NC$ bajo; uno cuyos vecinos son
casi todos del mismo tipo tiene $NC$ alto (centralidad estructural).

\textbf{Uso típico}: \emph{penaliza} paths que atraviesan hubs
naturales (ej.\ un DC con 5000 usuarios conectados) porque su alto
score en otras dimensiones se explica por topología, no por anomalía
real.

\section{Componente 4: Temporal Novelty}
\label{sec:scoring-tn}

Qué tan recientemente apareció la entidad en la ventana observada.

\begin{equation*}
TN(v) = \frac{\tau_{\text{first}}(v) - \tau_{\min}}{\tau_{\max} - \tau_{\min}}
\qquad C_4(p) = \frac{1}{|p|} \sum_{v \in p} TN(v)
\end{equation*}

$TN(v) \approx 0$: el nodo existe desde el inicio de la ventana.
$TN(v) \approx 1$: el nodo apareció al final.

\textbf{Uso típico}: separa \emph{centralidad estructural} de
\emph{novedad genuina}. Un DC corporativo tiene $TN = 0$ (siempre
estuvo); una IP externa recién observada tiene $TN \approx 1$.

\section{Componente 5: GNN Threat}
\label{sec:scoring-gnn}

Score de amenaza por path según clasificador GNN entrenado
(ver capítulo de Arquitectura, Capa Core). Cuando no hay modelo
cargado, devuelve 0.0 para todo path — caso común en entornos
limpios.

Clases de amenaza soportadas (\code{ThreatClass}):
\code{Benign}, \code{Exfiltration}, \code{C2Beacon},
\code{LateralMovement}, \code{PrivilegeEscalation}. La probabilidad
del argmax se usa como score. Peso por defecto 0.0 (opt-in).

\section{Explicación}

El endpoint \code{GET /explain_score} y la tool
\code{scoring.explain_score} devuelven el desglose:

\begin{lstlisting}[basicstyle=\ttfamily\small]
{
  "path_id": "...",
  "composite": 0.73,
  "components": {
    "entity_rarity":              0.82,
    "edge_rarity":                0.91,
    "neighborhood_concentration": 0.45,
    "temporal_novelty":           0.64,
    "gnn_threat":                 0.00
  },
  "weights": { "w1": 0.25, "w2": 0.30, "w3": 0.25, "w4": 0.20, "w5": 0.0 },
  "narrative": "entity rare (0.82), transition rare (0.91), ..."
}
\end{lstlisting}

\code{narrative} es texto humano generado en
\filepath{core/graph-engine/src/analytics.rs::build_narrative}. No
substituye el desglose numérico; lo acompaña.

\section{Invariantes}

\begin{itemize}
    \item Cada $C_i(p) \in [0, 1]$.
    \item $S(p) \in [0, 1]$ estrictamente tras renormalización.
    \item Path vacío devuelve todos los componentes en 0.0 sin
    panic.
    \item Si el grafo no está finalizado, el scorer cortocircuita a
    0.0 (precondición, no error).
    \item Todos los pesos en 0 $\Rightarrow$ $S(p) = 0$ para
    cualquier $p$.
\end{itemize}
